% !TEX root =../main.tex

\section{Conclusiones}
Las redes de comunicaciones complejas pueden conformarse por múltiples instancias de un
mismo tipo de dispositivo, formando topologías que, a simple vista, podrían parecer
redundantes o ineficientes. Cuando una topología contiene conexiones redundantes (por
ejemplo, una malla donde cuatro dispositivos intermedios poseen tres conexiones entre
sí), existe el riesgo crítico de generar bucles de red. Estos bucles, al circular
paquetes de manera infinita, pueden provocar la caída del sistema en poco tiempo,
independientemente de si se trata de un entorno de laboratorio o de una infraestructura
corporativa con un flujo considerable de información.

Para mitigar este riesgo, se implementa el protocolo STP (Spanning Tree Protocol), el
cual opera de manera nativa en los switches. A través del escenario de prueba
planteado, utilizando cuatro switches y cuatro terminales, fue posible observar el
funcionamiento de este protocolo en un entorno real. Tras completar la configuración,
se verificó la efectividad del protocolo mediante la elección de un Root Bridge y la
identificación mutua entre conmutadores para habilitar la interacción. La estabilidad
de la red se validó mediante pruebas ICMP (ping), confirmando que el tráfico recorre la
topología redundante de manera exitosa y sin colisiones.

Finalmente, es imperativo reconocer que una red redundante se diseña para garantizar la
alta disponibilidad y la tolerancia a fallos. Para comprender la verdadera utilidad de
STP, se evidenció la convergencia de la red tras un evento inesperado: la desconexión
del Root Bridge. Este incidente forzó una reestructuración automática de los canales
redundantes, demostrando que el protocolo es capaz de restablecer la conectividad y
encontrar rutas alternativas en un tiempo mínimo, asegurando así la continuidad
operativa del sistema.

A diferencia del estándar STP (802.1D), cuyo proceso de convergencia depende de estados
temporizados que pueden resultar lentos, el protocolo RSTP (Rapid Spanning Tree
Protocol) introduce mecanismos de sincronización más ágiles. Mientras que el STP
convencional debe esperar tiempos de escucha y aprendizaje, el RSTP utiliza un sistema
de propuesta y acuerdo (handshake) entre switches vecinos para transicionar casi
instantáneamente al estado de reenvío. Esta evolución permite que, en teoría, la
topología se adapte a los cambios con una interrupción mínima, optimizando la
disponibilidad de la infraestructura mediante la asignación de roles de puerto alternos
y de respaldo.

No obstante, durante el desarrollo de las pruebas en el escenario planteado, se
presentó un comportamiento anómalo al evaluar la recuperación bajo el estándar RSTP.
Tras forzar la caída del Root Bridge para observar la reestructuración de la red, se
evidenció que uno de los dispositivos no respondió satisfactoriamente al proceso de
sincronización. Este fallo en la comunicación técnica provocó que el equipo realizara
un fallback automático hacia el protocolo STP original por razones de
retrocompatibilidad. Como resultado, la topología tardó más tiempo en estabilizarse que
en la prueba inicial de STP, demostrando que factores externos o inconsistencias en la
negociación del hardware pueden degradar el rendimiento de protocolos que, en
condiciones ideales, deberían ser superiores.


A partir de este análisis, es posible concluir que una topología redundante ofrece la
ventaja fundamental de la alta disponibilidad, permitiendo que la red encuentre rutas
alternativas ante eventos inesperados. Sin embargo, su principal desventaja radica en
la complejidad técnica que esto conlleva; la dependencia de protocolos de control
introduce variables críticas donde un error de negociación puede generar tiempos de
convergencia impredecibles. Esto resalta que la robustez de una red no solo depende de
tener múltiples conexiones, sino de una configuración precisa y una compatibilidad
total entre los equipos para garantizar que la recuperación sea tan eficiente como la
teoría lo predice.

